\section{AIの社会的論点と知識集約社会}

\subsection{令和元年度再試験 問15}
\label{q:r01-ret-apt-15}

\noindent\textbf{原典：} sources/R01\_retest/aptitude.pdf、PDF物理16頁。

\textbf{公式正答：} 1（single）。

\paragraph{問題の要旨} AIの感情誘導・価値評価・データ保護・技術格差を2017年の六論点から判定する。

\Needspace{4\baselineskip}\paragraph{確認した条件・記述（独自要約）}\mbox{}\par

\noindent\textbf{原典確認：} 問題幹と各記述の主体・否定・全称・表の5候補を保存原典画像で確認。\par

\noindent\textbf{制度時点：} R01年度の原典表現を保持。現行資料の条番号や仕様を過去へ補わない。\par

\noindent\textbf{各記述の判断：} ア：○：本人が気付かない感情等の誘導は倫理課題である。；イ：○：生産性と創作価値の多様性を考える必要がある。；ウ：○：プライバシー保護と利活用の過度な萎縮防止を両立させる。；エ：○：AIリテラシーの差によるデバイドが起こり得る。\par

\noindent\textbf{前提：} 2017年報告書は倫理・法・経済・教育・社会・研究開発の六観点で便益と課題を検討する。社会の価値と機械の精度は別であり、利用できる能力の格差も課題となる。\par

\paragraph{問われる核心} AIの感情誘導・価値評価・データ保護・技術格差を2017年の六論点から判定する。

\paragraph{必要知識} 2017年報告書は倫理・法・経済・教育・社会・研究開発の六観点で便益と課題を検討する。社会の価値と機械の精度は別であり、利用できる能力の格差も課題となる。

\paragraph{考え方・正答を選ぶ理由（査読待ち初稿）} 四記述は報告書の倫理・法・社会の検討に対応する。不適切0個、1である。生産性向上は健全な利活用への期待であり、どの導入でも必ず達成されるという定理ではない。格差の可能性を、利用しない人への強制と取り違えない。

\Needspace{4\baselineskip}\paragraph{選択肢の検討（査読待ち初稿）}\mbox{}\par

\noindent\textbf{1：適切} 0。不適切なものは0個であり、一致する。\par

\noindent\textbf{2：不適切} 1。不適切なものは0個であり、この個数とは一致しない。\par

\noindent\textbf{3：不適切} 2。不適切なものは0個であり、この個数とは一致しない。\par

\noindent\textbf{4：不適切} 3。不適切なものは0個であり、この個数とは一致しない。\par

\noindent\textbf{5：不適切} 4。不適切なものは0個であり、この個数とは一致しない。\par

\paragraph{再利用する考え方} AIの技術能力、制度、価値判断、利用の機会を別軸で考える。

\paragraph{共通理論} \hyperref[theory:ai:society]{第\ref*{theory:ai:society}節}

\par 参考：\cite{ai_society2017}。

\paragraph{関連問題} \hyperref[q:r03-reg-apt-06]{R03-REG-APT-06}

\par 難易度 3（校準初稿）、優先度未評価。

\subsection{令和7年度 問7}
\label{q:r07-reg-apt-07}

\noindent\textbf{原典：} sources/R07/aptitude.pdf、PDF物理7頁。

\textbf{公式正答：} 1（single）。

\paragraph{問題の要旨} 2024年白書の知識集約社会と生成AIの普及を四つの語句へ対応させる。

\Needspace{4\baselineskip}\paragraph{確認した条件・記述（独自要約）}\mbox{}\par

\noindent\textbf{原典確認：} 問題幹と各記述の主体・否定・全称・表の5候補を保存原典画像で確認。\par

\noindent\textbf{制度時点：} R07年度の原典表現を保持。現行資料の条番号や仕様を過去へ補わない。\par

\noindent\textbf{各記述の判断：} ア：知識；イ：高付加価値化；ウ：医療；エ：対話\par

\noindent\textbf{前提：} 白書は知識を価値創造の中心とする社会、製品やサービスの高付加価値化、医療への影響、対話型インターフェースの普及を説明する。\par

\paragraph{問われる核心} 2024年白書の知識集約社会と生成AIの普及を四つの語句へ対応させる。

\paragraph{必要知識} 白書は知識を価値創造の中心とする社会、製品やサービスの高付加価値化、医療への影響、対話型インターフェースの普及を説明する。

\paragraph{考え方・正答を選ぶ理由（査読待ち初稿）} 2024年白書の対象本文は知識集約型、高付加価値化、医療、対話型生成AIと記している。四語が全て合う1である。AIの利活用の一般的な倫理課題が存在することから、本文の医療という産業分野を倫理へ置き換えない。

\Needspace{4\baselineskip}\paragraph{選択肢の検討（査読待ち初稿）}\mbox{}\par

\noindent\textbf{1：適切} 知識・高付加価値化・医療・対話。四空欄が白書の対象記述に対応する。\par

\noindent\textbf{2：不適切} 情報・標準規定化・倫理・対話。知識集約・高付加価値化・医療の三語が違う。\par

\noindent\textbf{3：不適切} 知識・標準規定化・医療・学習。高付加価値化と対話型の二語が違う。\par

\noindent\textbf{4：不適切} 情報・標準規定化・医療・対話。知識集約と高付加価値化の二語が違う。\par

\noindent\textbf{5：不適切} 知識・高付加価値化・倫理・学習。医療と対話型の二語が違う。\par

\paragraph{再利用する考え方} 政策資料の時点と対象箇所を固定し、一般的な連想で空欄を補わない。

\paragraph{共通理論} \hyperref[theory:ai:knowledge_society]{第\ref*{theory:ai:knowledge_society}節}

\par 参考：\cite{mext_ai_whitepaper2024}。

\par 難易度 3（校準初稿）、優先度未評価。

\section{AI利活用の十原則}

\subsection{令和3年度 問6}
\label{q:r03-reg-apt-06}

\noindent\textbf{原典：} sources/R03/aptitude.pdf、PDF物理6頁。

\textbf{公式正答：} 1（single）。

\paragraph{問題の要旨} 2019年AI利活用ガイドラインの十原則に含まれないものを数える。

\Needspace{4\baselineskip}\paragraph{確認した条件・記述（独自要約）}\mbox{}\par

\noindent\textbf{原典確認：} 問題幹と各記述の主体・否定・全称・表の5候補を保存原典画像で確認。\par

\noindent\textbf{制度時点：} R03年度の原典表現を保持。現行資料の条番号や仕様を過去へ補わない。\par

\noindent\textbf{各記述の判断：} ア：適正利用；イ：適正学習；ウ：連携；エ：安全；オ：セキュリティ；カ：プライバシー；キ：尊厳・自律；ク：公平性；ケ：透明性；コ：アカウンタビリティ\par

\noindent\textbf{前提：} 利用者にはサービス提供者・業務利用者・消費者が含まれる。適正利用、適正学習、連携、安全、セキュリティ、プライバシー、尊厳自律、公平性、透明性、説明責任の十原則を整理する。\par

\paragraph{問われる核心} 2019年AI利活用ガイドラインの十原則に含まれないものを数える。

\paragraph{必要知識} 利用者にはサービス提供者・業務利用者・消費者が含まれる。適正利用、適正学習、連携、安全、セキュリティ、プライバシー、尊厳自律、公平性、透明性、説明責任の十原則を整理する。

\paragraph{考え方・正答を選ぶ理由（査読待ち初稿）} 十記述は2019年ガイドライン9--10頁の十原則に全て対応する。含まれないもの0個、1である。ただしガイドライン5--6頁は利用の目的・文脈に応じた原則の選定を説明しており、全ての利用者へ全措置を同じ形で要求する法律ではない。

\Needspace{4\baselineskip}\paragraph{選択肢の検討（査読待ち初稿）}\mbox{}\par

\noindent\textbf{1：適切} 0。含まれないなものは0個であり、一致する。\par

\noindent\textbf{2：不適切} 1。含まれないなものは0個であり、この個数とは一致しない。\par

\noindent\textbf{3：不適切} 2。含まれないなものは0個であり、この個数とは一致しない。\par

\noindent\textbf{4：不適切} 3。含まれないなものは0個であり、この個数とは一致しない。\par

\noindent\textbf{5：不適切} 4。含まれないなものは0個であり、この個数とは一致しない。\par

\paragraph{再利用する考え方} 原則名の列挙と、各主体へ必要な具体策の選定を区別する。

\paragraph{共通理論} \hyperref[theory:ai:principles]{第\ref*{theory:ai:principles}節}

\par 参考：\cite{ai_use2019}。

\paragraph{関連問題} \hyperref[q:r01-ret-apt-15]{R01-RET-APT-15}

\par 難易度 3（校準初稿）、優先度未評価。
